% Vorlage fuer Abschlussarbeiten am Fachbereich Mathematik der TU Darmstadt.
% Geeignet fuer Bachelorarbeiten, Masterarbeiten und Diplomarbeiten.

%\documentclass[12pt,a4paper,twoside]{report}
\documentclass[12pt,a4paper,leqno]{article}
\usepackage[paper=a4paper,left=25mm,right=25mm,top=25mm,bottom=25mm]{geometry}
\usepackage[ngerman]{babel}
\usepackage[utf8]{inputenc}
\usepackage[T1]{fontenc}

\usepackage{enumerate}
\usepackage{amsmath, amssymb, amsthm, dsfont}
\usepackage{mathtools}
\usepackage{tikz}
\usepackage{tikz-cd}
\usetikzlibrary{quotes,babel,angles}
%\usepackage{hyperref}
\usepackage[pdfauthor={Fabian Gabel},%
pdftitle={Topologie-Seminar - Reflektionen \& Coreflektionen},%
%pagebackref=true,%
%pdftex
]{hyperref}
\usepackage{lmodern}

\parindent 0pt

\linespread{1.2}

% Inhaltsverzeichnis erstellen.

../ausarbeitung/befehle.tex

\begin{document}

\pagestyle{empty}
\pagenumbering{roman}
\hypersetup{pageanchor=false}
\begin{titlepage}
  \begin{center}
    \vspace{1cm}
    \includegraphics[width=0.5\linewidth]{TU_Darmstadt_Logo.pdf}
    \vspace{12pt}
    
    \large{Fachbereich Mathematik}
    \vspace{2cm}
    
    \large{Topologie-Seminar im Sommersemester 2017}
    \vspace{2cm}

    \huge{Reflektionen \& Coreflektionen}
    
    \vspace*{2cm}    
    
		\large
                \href{mailto:gabel@mathematik.tu-darmstadt.de}{Fabian Gabel}
    \vspace*{.5cm}

    01.06.2017 \\
    \vspace*{1cm}

    Veranstalter: Dr. rer. nat. Ren\'e Bartsch

    \vspace*{.5cm}

    %Zweiter Gutachter: Name des zweiten Gutachters\\[2ex]
    \tiny{Version vom \today}
  \end{center}
\end{titlepage}
\vspace*{\fill}

\pagestyle{plain}

\hypersetup{pageanchor=true}
\pagenumbering{arabic}

\tableofcontents

%\pagestyle{headings}

\section*{Einleitung}

So wie innerhalb einer Kategorie einzelne Objekte zueinander über Morphismen in Beziehung stehen, lassen sich auch Kategorien zueinander in Beziehung setzen. Statt Morphismen spricht man in diesem Zusammenhang von Funktoren. Man könnte den Blickwinkel beziehen die Kategorie der Kategorien zu betrachten, in welcher die Objekte durch Kategorien und die Morphismen durch Funktoren repräsentiert werden, um zu erkennen, dass letztlich wieder die Isomorphiefrage im Raum steht.

Isomorphie von Kategorien ist nur allzu oft nicht anzutreffen und dies macht in gewisser Weise auch den Reiz unterschiedlicher Kategorien aus. In einer ersten Abschwächung der Isomorphie wird man hoffen eine Äquivalenz von Kategorien zu finden. Ist auch das noch zu viel verlangt, so möchte man dennoch versuchen zwei Funktoren, die in entgegengesetzter Richtung zwischen zwei Kategorien operieren, mit den jeweiligen Identitätsfunktoren auf \emph{natürliche} Art in Verbindung zu bringen. Man kommt zum Begriff der Adjunktion, um den es sich in dieser Ausarbeitung drehen wird. Insbesondere werden wir uns mit Inklusionsfunktoren in topologischen Konstrukten beschäftigen und einige der zwischen topologischen Kategorien bestehenden Relationen kategorientheoretisch beschreiben.

Die folgende Ausarbeitung beschränkt sich bis auf ein paar Ausnahmen darauf Resultate aus dem Buch \cite{preuss} zusammenzufassen und hierbei größtenteils auf Beweise zu verzichten. 
Dies sollte keinesfalls auf die Faulheit des Erstellers zurückgeführt werden, sondern eher als Empfehlung verstanden werden, entsprechende Passagen im besagten Buche nachzulesen, da die Beweise dort aus Sicht des Erstellers bereits in einer verständlichen ausführlichen Form vorliegen und eine Aufnahme dieser in die vorliegende Ausarbeitung nicht zu derer Übersichtlichkeit beigetragen hätte.
In diesem Sinne ist diese Ausarbeitung eher als Wegbeschreibung durch das zweite Kapitel aufzufassen, durchaus aber auch als ein Appetithappen.

\section{Kategorientheoretische Grundlagen (Fortsetzung)}

In diesem Abschnitt füllen wir das Vokabelheft mathematischer Definitionen mit weiteren Begriffen aus der Kategorientheorie.

\subsection{Funktoren, universelle Morphismen und natürliche Transformationen}

Wie schon in der Einleitung angekündigt, werden Funktoren in diesem Teil des Seminars eine prominente Rolle spielen, denn wir wollen unterschiedliche Kategorien zueinander in Relation setzen und Abbildungen, die dies auf \emph{natürliche} Weise schaffen, nennen wir Funktoren.

\begin{defn}[Covarianter Funktor]
  Seien $\C$ und $\De$ Kategorien und $\F_1 \colon |\C| \to |\De|$ und $\F_2 \colon \Mor_\C\to \Mor_\De$. 
  Dann nennen wir das Quadrupel $\F = (\C, \De, \F_1, \F_2)$ einen (\emph{covarianten}) \emph{Funktor} von $\C$ nach $\De$, falls folgende Bedingungen erfüllt sind:
  \begin{enumerate}[F1)]
    \item $f \in [A, B]_\C$ impliziert $\F(f) \in [\F(A), \F(B)]_\De$.
    \item $\F(f \circ g) = \F(f) \circ \F(g)$, falls $f \circ g$ definiert ist.
    \item $\F(1_A) = 1_{\F(A)}$ für alle $A \in |\C|$.
  \end{enumerate}
  Abkürzend schreiben wir im Folgenden auch $\F \colon \C \to \De$.
\end{defn}

Aufgrund der Eigenschaften F2) und F3) bezeichnet man Funktoren manchmal auch als Homomorphismen von Morphismen.

Ein \emph{contravarianter Funktor} $\F \colon \C \to \De$ ist gerate ein covarianter Funktor von $\C^* \to \De$, was gerade bedeutet, dass folgende modifizierten Bedingungen gelten:
\begin{enumerate}[F1')]
  \item $f \in [A, B]_\C$ impliziert $\F(f) \in [\F(B), \F(A)]_\De$.
  \item $\F(f \circ g) = \F(g) \circ \F(f)$, falls $f \circ g$ existiert.
\end{enumerate}

Wir beleben den Begriff des Funktors nun, indem wir bekannte Sachverhalte kategorientheoretisch unter die Lupe nehmen.

\begin{ex}
  \label{ex:Example}
\begin{enumerate}[a)]
    \item Konstanter Funktor: Sind $\C$ und $\De$ Kategorien und $X \in |\De|$, so lässt sich für alle $A \in |\C|$ und alle $f \in \Mor_\C$ durch $\F(A) = X$ und $\F(f) = 1_X$ ein Funktor definieren, welcher sowohl covariant als auch contravariant ist.
    \item Vergissfunktor: Ist $\C$ ein (topologisches) Konstrukt, so lässt sich ein Funktor $\F \to \Set$ definieren durch $\F((X, \xi)) = X$ und $\F(f) = f$.
    \item Dualisierender Funktor $\Delta_\C$: Es lässt sich ein Funktor $\F \colon \C \to \C^*$ definieren durch $\F(X) = X$ und $F(f) = f^*$. Dieser ist natürlich contravariant.
    \item Dualer Funktor: Ist $\F \colon \C \to \De$ ein Funktor, so erhält man den zugehörigen dualen Funtor als $\Delta_{\De} \circ F \circ \Delta_{\C^*}$
    \item Inklusionsfunktor: Sei $\C$ eine Kategorie und $\A$ eine \emph{Unterkategorie}, also eine Kategorie für die gilt
      \begin{enumerate}
        \item $|\A| \subset |\C|$,
        \item $[A, B]_\A \subset [A,B]_\C$ für alle $(A, B) \in |\A| \times |\A|$,
        \item Die Komposition von Morphismen in $\A$ stimmt mit der Komposition in $\C$ überein und die der Identitätsmorphismus ist auch derselbe.
      \end{enumerate}
      Gilt zu alledem sogar $[A, B]_\A = [A, B]_\C$, so bezeichnen wir $\A$ als \emph{volle} Unterkategorie.
    \item Identitätsfunktor $\I_\C$: Dieser Funktor $\F \colon \C \to \C$ wird durch $\F(X) = X$ und $\F(f) = f$ definiert.
  \end{enumerate}
\end{ex}

Wir sehen also: Unsere Welt ist voller Funktoren! Um den Kreis der Verdächtigen, die wir später noch genauer unter die Lupe nehmen wollen, etwas einzuschränken, führen wir einen speziellen Typ von Morphismen ein, der es uns im weiteren Verlauf gestatten wird, verdächtige kategorientheoretische Sachverhalte als funktorielle Zusammenhänge zu entlarven.

\begin{defn}[Universelle Abbildung]
  Seien $\A$ und $\B$ Kategorien, $\F \colon \A \to \B$ ein Funktor und $B \in |\B|$.
  Ein Paar $(u, A)$ mit $A \in |\A|$ und $u \colon B \to \F(A)$ heißt \emph{universelle Abbildung für $B$ bezüglich $\F$}, falls für alle $A' \in |\A|$ und alle $f \colon B \to \F(A')$ genau ein $\A$-Morphismus $\overline f \colon A \to A'$ existiert, sodass das Diagramm
  $$
  \begin{tikzcd}[row sep=2.5em]
    B \arrow[rd,"u"] \arrow[rr, "f"] &&\F(A') \\
    &\F(A) \arrow[ru, "\F(\overline{f})"]
  \end{tikzcd}
  $$
  kommutiert.
  Entsprechend bezeichnet man ein Paar $(A,u)$ mit $A \in |\A|$ und $u \colon \F(A) \to B$ als \emph{co-universelle Abbildung für $B$ bezüglich $\F$}, falls $(u^*, A)$ eine universelle Abbildung für $B$ bezüglich des dualen Funktors $\F^* \colon \A^* \to \B^*$ ist.
  Dies bedeutet, dass für alle $A' \in |\A|$ und jeden $\B$-Morphismus $f \colon \F(A') \to B$ ein eindeutiger $\A$-Morphismus $\overline f \colon A' \to A$ existiert, so dass das Diagramm
  $$
  \begin{tikzcd}[row sep=2.5em]
    B  &&  \arrow[ll, "f"] \F(A')\arrow[ld, "\F(\overline{f})"]  \\
    & \arrow[lu,"u"]\F(A) & 
  \end{tikzcd}
  $$
  kommutiert.
\end{defn}

Man beachte, dass der Morphismus $\F(\overline f)$ keineswegs eindeutig festgelegt sein muss. Dies ist im Allgemeinen erst dann der Fall, wenn $\F$ ein \emph{treuer} Funktor ist.

Im folgenden Lemma beschreiben wir das Verhalten (co-)universeller Abbildungen unter Verknüpfung mit Isomorphismen.

\begin{lem}
  \label{lem:universalCircIso}
  Seien $\A$ und $\B$ Kategorien, $\F \colon \A \to \B$ ein Funktor und $B \in |\B|$ und $(u,A)$ eine universelle Abbildung für $B$ bezüglich $\F$.
  Sei nun $v \colon A \to \underline A$ ein $\A$-Isomorphismus, dann ist auch $(\F(v) \circ u, \underline A )$ eine universelle Abbildung für $B$ bezüglich $\F$.

  Ist $(A, u)$ eine couniverselle Abbildung für $B$ bezüglich $\F$, so ist auch $(\underline A, u \circ \F(v^{-1}))$ eine couniverselle Abbildung für $B$ bezüglich $\F$.
\end{lem}

\begin{proof}
  Sei $f \colon B \to \F(A')$ ein $\B$-Morphismus. 
  So existiert aufgrund der Eigenschaften von $u$ genau ein $\A$-Morphismus $\overline f \colon A \to A'$ mit $f = \F(\overline f) \circ u$.
  Aufgrund der Eindeutigkeit von $\overline f$ existiert somit auch genau ein $g \coloneqq v^{-1} \circ \overline f \colon \underline A \to A'$, sodass das Diagramm   
  $$
  \begin{tikzcd}[row sep=2.5em]
    B \arrow[rd,"u"] \arrow[rr, "f"] &&\F(A') \\
    &\F(A) \arrow[ru, "\F(\overline{f})"'] \arrow[rd,"\F(v)"] \\
    &&\F(\underline A) \arrow[uu,"\F(g)"']
  \end{tikzcd}
  $$
  mitsamt seiner Unterdiagramme kommutiert.
  Über ein analoges Argument zeigt man, dass im Falle einer couniversellen Abbildung $(A,u)$ das Diagramm
  $$
  \begin{tikzcd}[row sep=2.5em]
    B  & & \F(A') \arrow[ll, "f"]\arrow[ld, "\F(\overline{f})"] \arrow[dd, "\F(g)"] \\
    & \F(A) \arrow[lu,"u"] \arrow[rd, "\F(v)"]  & \\
    & & \F(\underline A) 
  \end{tikzcd}
  $$
  kommutiert. Hierbei ist $g \coloneqq v \circ \overline f$.
\end{proof}

Wir halten nun gewissermaßen die Umkehrung des vorangehenden Lemmas fest, nämlich, dass universelle Abbildungen bereits eindeutig bis auf Isomorphie sind.

\begin{prop}[\cite{preuss}, 2.1.6]
  Seien $(u,A)$ und $(u',A')$ universelle Abbildungen für $B \in |\B|$ bezüglich $\F \colon \A \to \B$.
  Dann existiert ein Isomorphismus $f \colon A \to \A'$, sodass das Diagramm
  $$
  \begin{tikzcd}[row sep=2.5em]
    B \arrow[rd,"u'"] \arrow[rr, "u"] &&\F(A) \\
    &\F(A') \arrow[ru, "\F(\overline{f})"]
  \end{tikzcd}
  $$
  kommutiert.
\end{prop}

Widmen wir unsere Aufmerksamkeit nun einem famosen Beispiel, um die eigeführten Prinzipien bei der Arbeit zu bestaunen.

\begin{ex}[Stone-\v{C}ech Kompaktifizierung]
       Sei $\Tych$ die Kategorie der Tychonoff-Räume, und $\CompHaus$ die Kategorie der kompakten Hausdorff-Räume.
      Ferner bezeichnen wir für $X \in |\Tych|$ mittels $\beta(X)$ seine Stone-\v{C}ech-Kompatifizierung und mit $e_X \colon X \to \beta(X)$ die entsprechende Einbettung. Dann ist nach dem Satz von Stone-\v{C}ech \cite[5.4.8]{bartsch} das Paar $(e_X, \beta(X))$ eine universelle Abbildung bezüglich des Inklusionsfunktors $\F_e \colon \CompHaus \to \Tych$.
      Ein entsprechend angepasstes kommutatives Diagramm liefert Gewissheit:
      $$
      \begin{tikzcd}[row sep=2.5em]
        X \arrow[rd,"e_X"] \arrow[rr, "f"] &&\F_e(Y) = Y \\
        &\F_e(\beta(X)) = \beta(X) \arrow[ru, "\F_e(\overline{f})"]
      \end{tikzcd}
      $$
      Hierbei sei $Y \in \CompHaus$ und $f \in [X, \F_e(Y)]_{\Tych}$.
\end{ex}

Funktoren für sich alleine sind ja schon interessante Kreaturen, aber wie lassen sich Beziehungen unterschiedlicher Funktoren zueinander beschreiben?
Stellt man sich die Frage, ob man nicht auch Funktoren als Objekte eine Kategorie auffassen kann, so kommt sofort die Frage auf, welches die richtigen Morphismen zwischen Funktoren sind. 

\begin{defn}[Natürliche Transformationen]
  Seien $\C$ und $\De$ Kategorien und $\F, \G \colon \C \to \De$ Funktoren.
  \begin{enumerate}[1)]
    \item Eine Familie $\eta = (\eta_A)_{A \in |\C|}$ mit $\eta_A \in [\F(A), \G(A)]_\De$ für alle $\A \in |\C|$ heißt \emph{natürliche Transformation}, falls für alle $(A,B) \in |\C| \times |\C|$ und alle $f \in [A,B]_\C$ das rechte Diagramm
      $$
      \begin{tikzcd}[row sep=2.5em]
        A\arrow[d,"f"] & & \F(A)\arrow[r,"\eta_A"]\arrow[d,"\F(f)"'] & \G(A)\arrow[d,"\G(f)"] \\
        B & & \F(B)\arrow[r,"\eta_B"] & \G(B)
      \end{tikzcd}
      $$
      kommutiert. Ist $\eta$ eine natürliche Transformation, schreiben wir im Folgenden auch kurz $\eta \colon \F \to \G$.
    \item Eine natürliche Transformation $\eta \colon \F \to \G$ heißt \emph{natürliche Äquivalenz}, falls für alle $A \in |\C|$ der Morphismus $\eta_A$ ein Isomorphismus ist.
    \item $\F$ und $\G$ heißen \emph{natürlich äquivalent}, oder kurz $\F \approx \G$, wenn eine natürliche Äquivalenz $\eta \colon \F \to \G$ existiert.
  \end{enumerate}
\end{defn}

Weiteres zur Kategorie der Funktoren lässt sich in \cite{brandenburg} nachlesen.

\subsection{Adjungierte Funktoren}

Ein Funktor $\F \colon \C \to \De$ ist natürlich nicht immer ein Isomorphismus, d.h. es existiert nicht unbedingt ein Funktor $\G \colon \De \to \De$ mit
$$
\F \circ \G = \I_\De \quad\text{und}\quad \G \circ \F = \I_\C.
$$

Sind Funktoren nicht invertierbar, können sie zumindest \emph{fast} invertierbar mittels einer Pseudoinversen $\G \colon \De \to \C$ sein.
Dies ist ein Funktor mit 
$$\F \circ \G \approx \I_\De \quad\text{und}\quad \G \circ \F \approx \I_\C.$$

Auch das ist nicht immer der Fall und folglich ist man bestrebt zu untersuchen, ob ein Funktor nicht zumindest \emph{fast fast} invertierbar ist, im Sinne dass ein weiterer Funktor $\G \colon \De \to \C$ existiert, sodass zumindest natürliche Transformationen
$$
\eta \colon \F \circ \G \to \I_\De \quad\text{und}\quad \xi \colon \G \circ \F \to \I_\C
$$
zur Stelle sind.
Um ebendiese Funktoren dreht es sich in diesem Unterabschnitt.

Zunächst wollen wir jedoch die im vorangehenden Kapitel neu eingeführten Begriffe \emph{natürliche Transformation} und \emph{universelle Abbildung} mit einander verheiraten.

\begin{thm}[\cite{preuss}, 2.1.12]
  Sei $\F \colon \A \to B$ ein Funktor mit der Eigenschaft, dass für alle $B \in |\B|$ eine universelle Abbildung $(u_B, A_B)$ bezüglich $F$ existiert.
  Dann existiert genau ein Funktor $\G \colon B \to \A$, sodass Folgendes gilt:
  \begin{enumerate}[(1)]
    \item $\G(B) = A_B$ für alle $B \in |\B|$.
    \item $u = (u_B) \colon \I_B \to \F \circ \G$ ist eine natürliche Transformation.
  \end{enumerate}
\end{thm}

Moralisch gesehen sollte der Morphismus-Part des entstehenden Funktors gegeben sein durch $\G(f) = \overline f$ und dies wird für die im Späteren betrachteten topologischen Konstrukte sogar nicht nur intuitiv richtig sein. 
Die saubere Realität sieht anders aus, wie der folgende Beweis zeigen soll.

\begin{proof}
  Die Aussage des Theorems beinhaltet bereits den Objekt-Teil des Funktors $\G$, den wir auch brav so definieren, d.h. für alle $B \in |B|$ setzen wir $\G(B) = A_B$.

  Betrachten wir nun das Diagramm
  $$
    \begin{tikzcd}[row sep=2.5em]
      \I_\B(B)  = B \arrow[r,"u_B"] \arrow[d,"f"'] & \F(A_B)     = \F(\G(B)) \arrow[d,"\F(\widetilde f)"] 
      \\ 
      \I_\B(B') = B'\arrow[r,"u_{B'}"]              & \F(A_{B'})  = \F(\G(B')) 
    \end{tikzcd}
    $$
    Hierbei sei $\widetilde f \coloneqq \overline{u_{B'} \circ f} \in [A_B, A_{B'}]_\A$ eindeutig festgelegt über die Eigenschaft der universellen Abbildung $u_B$.
    Wir widerstehen der Versuchung nicht und setzen einfach $\G(f) = \widetilde f$. 
    Tatsächlich haben wir aufgrund der Eindeutigkeitseigenschaft universeller Abbildungen auch gar keine andere Wahl.
    Hiermit ist schonmal folgendes klar: Ist $\G$ ein Funktor, so gibt es keinen weiteren Funktor, der denselben Job erfüllt. 
    Da obiges Diagramm kommutiert, erhalten wir, gesetzt dem Falle, dass $\G$ tatsächlich ein Funktor ist, zudem die im Theorem postulierte Eigenschaft (2).

    Seien also $f,g \in  \Mor_\B$, so erhalten wir das folgende Diagramm:
    $$
        \begin{tikzcd}[column sep=3.5em,row sep=2.5em]
          B \arrow[r,"u_B"] \arrow[d,"f"'] \arrow[dd,"g \circ f"',bend right=60]    & \F(A_B)    \arrow[d,"\F(\widetilde f)"] \arrow[dd, bend left=60, "\F(\widetilde g \circ \widetilde f)"] 
      \\ 
      B'\arrow[r,"u_{B'}"] \arrow[d,"g"'] & \F(A_{B'}) \arrow[d ,"\F(\widetilde g)"] \\
      B''\arrow[r,"u_{B''}"]              & \F(A_{B''}) 
    \end{tikzcd}
    $$
    Dieses Diagramm kommutiert, da $\F$ ein Funktor ist.
    Insbesondere gilt also $\F(\widetilde g) \circ \F(\widetilde f) = \F(\widetilde g \circ \tilde f)$. 
    Andererseits liefert die Definition von $u_B$ auch die Existenz einer Abbildung $\widetilde{g \circ f} = \overline{u_{B''} \circ g \circ f}$ mit 
    $$
    \F(\widetilde{g \circ f}) \circ u_B = g \circ f .
    $$
    Die Eindeutigkeitseigenschaft universeller Abbildungen liefert damit sofort $\widetilde{g \circ f} = \widetilde g \circ \widetilde f$ und folglich auch $\G(g \circ f) = \G(g) \circ \G(f)$.
    Analoge Haarspalterei liefert $\G(1_B) = 1_{A_B}$.  
\end{proof}

So wie ein Ingenieur gerne auf Knöpfe einer unbekannten Machine drückt, so setzt ein Mathematiker gerne mehr oder minder bekannte Morhpismen in neuartige Funktoren ein. Genauer gesagt, interessieren wir uns dafür, welche neuen Morphismen wir erhalten, wenn wir die universellen Abbildungen $(u_B, A_B)$ in den soeben konstruierten Funktor $\G$ einsetzen.

\begin{kor}[\cite{preuss}, 2.1.12]
  Es existiert genau eine natürliche Transformation $v = (v_A) \colon \G \circ F \to \I_A$, sodass das Folgende gilt:
  \begin{enumerate}[(a)]
    \item $\F(v_A) \circ u_{\F(A)} = \1_{\F(A)}$ für alle $A \in |\A|$.
    \item $v_{\G(B)} \circ \G(u_B) = \1_{\G(B)}$ für alle $B \in |\B|$.
  \end{enumerate}
\end{kor}

Wir machen nun den vorangehenden Satz und das zugehörige Korollar zur Definition und untersuchen die dadurch entstehenden Objekte.

\begin{defn}[Linksadjungierter Funktor]
  Sind $\F \colon \A \to \B$ und $\G \colon \B \to \A$ Funktoren und $u = (u_B) \colon \I_\B \to \F \circ \G$ sowie $v = (v_A) \colon \G \circ \F \to \I_\A$ natürliche Transformationen mit den Eigenschaften
  \begin{enumerate}[(1)]
    \item $\F(v_A) \circ u_{\F(A)} = \1_{\F(A)}$ für alle $A \in |\A|$ und
    \item $v_{\G(B)} \circ \G(u_B) = \1_{\G(B)}$ für alle $B \in |\B|$,
  \end{enumerate}
  so nennen wir $\G$ den \emph{zu $\F$ linksadjungierten Funktor} und analog nennen wir $\F$ den \emph{zu $\G$ rechtsadjungierten Funktor}.
  Das Paar $(\G, \F)$ nennen wir ein \emph{Paar adjungierter Funktoren}.
\end{defn}

Haben wir etwa zu viel versprochen?
Dem aufmerksame Leser wird sich wohl kaum ein \emph{Hey, wo sind meine universellen Abbildungen?} verkneinfen können.
Keine Angst.
Es stellt sich heraus, dass sich universelle Abbildungen gerne in Rudeln aufhalten, um natürliche Tranformationen zu bilden.

\begin{thm}[\cite{preuss}, 2.1.15]
  Ist $\G \colon \B \to \A$ ein zu $\F \colon \A \to \B$ linksadjungierter Funktor und $u = (u_B) \colon \I_\B \to \F \circ \G$ eine zugehörige natürliche Transformation, dann ist für alle $B \in |B|$ das Paar $(u_B, \G(B))$ eine universelle Abbildung bezüglich $\F$.
\end{thm}

\begin{bem}[Adjungierte Situation]
  Namen sind bekanntlich Schall und Rauch und in der Kategorientheorie nur dann zu gebrauchen, wenn die mit dem Namen assozierten Gebilde auf eine gutartige Weise eindeutig bestimmt sind.
  Dies ist bei adjungierten Funktoren samt Gefolgschaft der Fall:
  \begin{enumerate}[(i)]
    \item Adjungierte Funktoren sind eindeutig bestimmt bis auf natürliche Äquivalenz
    \item Die in der Definition von adjungierten Funktoren auftretenden natürlichen Transformationen sind eindeutig bis auf natürliche Äquivalenz
  \end{enumerate}
  Es macht somit Sinn einem paar adjungierter Funktoren natürliche Transformationen zuzuordnen. Ein so entstehendes Tupel $(\G, \F, u, v)$ wollen wir \emph{adjungierte Situation} taufen.
\end{bem}

\begin{ex}[Stone-\v{C}ech Kompaktifizierung]: Der Inklusionsfunktor $\F_e$ aus Beispiel \ref{ex:Example} besitzt also eine Linksadjungierte $\beta \colon \Tych \to \CompHaus$, die jedem Tychonoff-Raum $X$ seine Stone-\v{C}ech-Kompaktifizierung zuordnet. Beim Versuch, Tychonoff-Räume zu kompaktifizieren wurde also implizit ein ganzer Funktor mitkonstruiert.
\end{ex}



\section{Reflektive und coreflektive Unterkategorien}

In diesem Abschnitt wollen wir uns speziell mit Inklusionsfunktoren und ihren Adjungierten beschäftigen. 
In freier Wildbahn treten Inklusionsfunktoren unter anderem bei der Betrachtung von Unterkategorien auf, was sie vor allem für uns so interessant macht.

\subsection{Allgemeine Definitionen}

In diesem Unterabschnitt erweitern wir wieder unser kategorientheoretisches Vokabelheft.

\begin{defn}[Reflektive Unterkategorie]
  \label{defn:reflektor}
  Sei $A$ eine Unterkategorie einer Kategorie $\C$ und $\F_e \colon \A \to \C$ der Inklusionsfunktor.
  Dann nennen wir $\A$ \emph{reflektiv} in $\C$ genau dann, wenn eine der folgenden äquivalenten Bedingungen erfüllt ist:
  \begin{enumerate}[(1)]
    \item $\F_e$ besitzt den linksadjungierten Funktor $\R$.
    \item Jedes $X \in |C|$ besitzt eine universelle Abbildung $(r_X, X_\A)$ bezüglich $\F_e$.
  \end{enumerate}
  Den Funktor $\R$ nennen wir dann einen \emph{Reflektor}, die Morphismen $r_X \colon X \to X_\A$ nennen wir Reflektionen von $X$ bezüglich $\A$.

  Durch Dualisierung erhalten wir einen weiteren Begriff.
  Wir nennen $\A$ \emph{coreflektiv} in $\C$, genau dann, wenn $\A^*$ reflektiv ist in $\C^*$. 
\end{defn}

\begin{defn}
  In der Situation von Definition \ref{defn:reflektor} nennen wir $\A$ \emph{epireflektiv/ extremal epireflektiv/ bireflektiv in} $\C$, falls $\A$ reflektiv in $\C$ ist und der für alle $X \in |\C|$ existierende $\C$-Morphismus $r_X \colon X \to X_\A$ ein Epimorphismus/ extremaler Epimorphismus/ Bimorphismus ist.
  Die Morphismen $r_X$ nennen wir \emph{Epireflektionen/ extremale Epireflektionen/ Bireflektionen}.
\end{defn}

\subsection[Reflektoren und Coreflektoren in topologischen Konstrukten]{Reflektoren und Coreflektoren\\ in topologischen Konstrukten}

Im nächsten Kapitel werden wir uns mit Unterkategorien von topologischen Konstrukten beschäftigen. Während eine Unterkategorie immer Anlass zu einem Inklusionsfunktor gibt, ist im Allgemeinen nicht klar, dass zu besagtem Inklusionsfunktor ein passender Reflektor oder Coreflektor existiert. 

\begin{kor}[\cite{preuss}, 2.2.6]
  Ein volles und unter Isomorphie abgeschlossenes Unterkonstrukt $\A$ eines topologischen Konstruktes $\C$ ist bireflektiv in $\C$ genau dann, wenn es abgeschlossen ist unter Bildung von Produkten und initialen Unterobjekten.
\end{kor}

Wir erweitern etwas unseren kategorientheoretischen Wortschatz:

\begin{defn}
  Wir nennen ein Objekt $S$ einer Kategorie $\C$ \emph{Separator}, falls für alle paarweise verschiedenen Morphismen $f,g \colon A \to B$ mit gleichem Definitions- und Wertebereich ein Morphismus $h \colon S \to A$ existiert mit der Eigenschaft, dass $f \circ h \neq g \circ h$.
\end{defn}

Wir halten fest, dass jedes Objekt $(X,\xi)$ eines topologischen Konstrukts $\C$ mit $X \neq \emptyset$ ein Separator ist. 
Denn für zwei paarweise verschiedene Morphismem $f,g \colon (Y, \eta) \to (Z, \theta)$ unterscheiden sich die zugrundeliegenden Mengenabbildungen $f$ und $g$ zumindest schonmal in einem Punkt $y \in Y$. Betrachten wir nun die konstante Abbildung $h \colon X \to Y, h(x) = y$, so ist diese aufgrund der Voraussetzung $X \neq \emptyset$ wohldefiniert und zudem ein Morphismus.
Damit folgt sofort die Behauptung.

Es stellt sich heraus, dass die bloße Existenz von Separatoren weitere Eigenschaften der Coreflektionen in folgender Weise freilegt.

\begin{thm}[\cite{preuss}, 2.2.9]
  Sei $S$ ein Separator einer Kategorie $\C$ und $\A$ eine coreflektive Unterkategorie von $\C$, die $S$ enthält.
  Dann ist $\A$ bereits epicoreflektiv.
\end{thm}

Wir wissen also, wann eine coreflektive Unterkategorie epicoreflektiv ist. 
Der folgende Satz geht nun einen Schritt weiter zu bicoreflektiven Unterkategorien.

\begin{thm}[\cite{preuss}, 2.2.10]
  Sei $\A$ eine epicoreflektive Unterkategorie von $\C$. 
  Ist $\A$ zusätzlich eine \emph{volle} Unterkategorie, so ist $\A$ bereits bicoreflektiv.
\end{thm}

\begin{bem}[\cite{preuss}, 2.2.11]
  Jedes coreflektive, volle und unter Isomorphie abgeschlossene Unterkonstrukt $\A$ eines topologischen Konstrukts $\C$ ist bicoreflektiv, falls $|\A|$ mindestens ein Element mit nicht leerer zugrunde liegender Menge enthält.

  Wir wollen im Folgenden zeigen, dass in diesem Fall die Coreflektionen eine sehr einfache Gestalt annehmen.  
  Für $(X, \xi) \in |\C|$ ist die entsprechende Coreflektion $c_X \colon (Y_\A, \eta_\A) \to (X, \xi)$ bijektiv.
  Nach \cite[1.2.2.7]{preuss} existiert eine $\C$-Struktur $\xi_\A$ auf $X$, sodass $c_X \colon (Y_\A, \eta_\A) \to (X, \xi_\A)$ ein Isomorphismus ist.
  Da $\A$ nach Voraussetzung abgeschlossen unter Isomorphismen ist, gilt $(X, \xi_\A)\in |\A|$.
  Zudem ist $\xi_\A$ die gröbste aller $\C$-Strukturen $\xi'$, für die einerseits $\xi' \leq \xi$ und andererseits $(X, \xi') \in \A$ gilt.

  Nach Lemma \ref{lem:universalCircIso} ist also auch $c_X \circ c_X^{-1} = 1_X \colon (X,\xi_A) \to (X,\xi)$ eine universelle Abbildung, denn da $(Y_A,\eta_A)$ und $(X,\xi_A)$ Elemente aus $|\A|$ sind, ist der $\C$-Morphismus $c_x$ insbesondere ein $\A$-Morphismus, da $\A$ eine volle Unterkategorie von $\C$ ist.

  Daher ist $((X,\xi_A), 1_X)$ die Coreflektion von $(X,\xi)$ bezüglich $\A$, man erhält also bis auf Isomorphie die Coreflektion eines  $\C$-Objekts $(X,\xi)$ bezüglich $\A$ durch eine Modifikation der $\C$-Struktur $\xi$ auf $X$.
\end{bem}

Wir schließen nun dieses Kapitel mit einem letzten Resultat zu allgemeinen Topologischen Konstrukten, welches eine Antwort auf die Frage liefert, wie sich initiale und finale Strukturen auf topologische Unterkonstrukte übertragen.

\begin{thm}[\cite{preuss}, 2.2.12]
  Sei $\A$ ein volles und unter Isomorphie abgeschlossenes Unterkonstrukt eines topologischen Konstrukts $\C$.
  Dann ist auch $\A$ topologisch, vorausgesetzt dass $\A$ bireflektiv oder bikoreflektiv in $\C$ ist.
  
  Ist $\A$ bireflektiv (bicoreflektiv) in $\C$, dann stimmen die initialen (finalen) Strukturen in $\A$ mit denen aus $\C$ überein, während die finalen (initialen) Strukturen in $\A$ aus den finalen (initialen) Strukturen in $\C$ entstehen, indem man den Bireflektor (Bicoreflektor) anwendet.
\end{thm}


\section{Konvergenzstrukturen und uniforme Konvergenzstrukturen}

In diesem letzten Abschnitt betrachten wir nun unterschiedliche Konvergenzstrukturen und uniforme Strukturen durch die kategorientheoretische Brille, mit dem Ziel diese untereinander in Beziehung zu setzen und die Verbindung von Konvergenzstrukturen und uniformen Strukturen herzustellen.


\subsection{Konvergenzstrukturen}

Zunächst einmal halten wir fest, welche Konvergenzstrukturen für uns interessant sein werden.

\begin{defn}[GKonv und seine Kinder]
  Die Kategorie $\GConv$ der verallgemeinerten Konvergenzräume mit stetigen Abbildungen setzt sich wie folgt zusammen:
  \begin{enumerate}[a)]
    \item Für jede Menge $X$ sei $\FF(X)$ die Menge aller Filter auf $X$. 
      Ein \emph{verallgemeinerter Konvergenzraum} ist ein Paar $(X,q)$, wobei $X$ eine Menge und $q \subset \FF(X) \times X$ eine Relation von Filtern und Punkten (gegen die sie \emph{konvergieren}) ist. 
      Zusätzlich sollen folgende Eigenschaften erfüllt sein:
      \begin{enumerate}[C1)]
        \item $(\dot x, x) \in q$ für alle $x \in X$; \emph{alle Einpunktfilter konvergieren gegen ihren Erzeuger}.
        \item $(\G, x) \in q$, falls $(F,x) \in q$ und $G \supset F$; \emph{Oberfilter konvergenter Filter, erben Grenzwerte}
      \end{enumerate}
    \item Eine Abbildung $f \colon (X,q) \to (X',q')$ zwischen verallgemeinerten Konvergenzräumen heißt \emph{stetig}, falls für alle $(\F,x) \in q$ auch $(f(\F), f(x)) \in q'$ gilt.
  \end{enumerate}
  Ein verallgemeinerter Konvergenzraum heißt
  \begin{enumerate}[a)]
    \item[c)] \emph{Kent Konvergenzraum}, falls folgende Bedingung erfüllt ist:
      \begin{enumerate}
        \item[C3)] $(\F \cap \dot x, x) \in q$, falls $(F, x) \in q$; \emph{Abgeschlossenheit bezüglich endlicher Durchschnitte mit Einpunktfiltern}.
      \end{enumerate}
    \item[d)] \emph{Limesraum}, falls folgende Bedingung erfüllt ist:
      \begin{enumerate}
        \item[C4)] $(\F \cap \G, x) \in q$, falls $(\F,x) \in q$ und $(\G,x) \in q$; \emph{Abgeschlossenheit bezüglich endlicher Durchschnitte}
      \end{enumerate}
    \item[e)] \emph{Pseudotopologischer Raum}, falls folgende Bedingung erfüllt ist:
      \begin{enumerate}
        \item[C5)] $(\F,x) \in q$, falls $(\U, x) \in q$ für alle Ultrafilter $\U \supset \F$.
      \end{enumerate}
    \item[f)] \emph{Prätopologischer Raum}, falls folgende Bedingung erfüllt ist:
      \begin{enumerate}
        \item[C6)] $(\U_q(x), x) \in q$ für alle $x \in X$, wobei $\U_q(x) \coloneqq \bigcap\{ \F \in F(x) \colon (\F, x) \in q\}$
      \end{enumerate}
  \end{enumerate}
  Ein prätopologischer Raum $(X,q)$ heißt 
  \begin{enumerate}[a)]
    \item[g)] \emph{topologischer Raum}, falls die folgende Bedingung erfüllt ist:
      \begin{enumerate}
        \item[C7)] Für alle $U \in \U_q(x)$ existiert ein $V \in \U_q(x)$ sodass $U \in \U_q(y)$ für alle $y \in V$ gilt.
      \end{enumerate}
  \end{enumerate}
\end{defn}

Die eben definierten Klassen definieren volle und unter Isomorphie abgeschlossene Unterkonstrukte von $\GConv$, welche wir im Folgenden mit $\Lim$, $\PsTop$, $\PrTop$ und $\TPrTop$ bezeichnen werden.

\begin{bem}[\cite{preuss}, 2.3.1.2]
Entsprechend der Definitionsreihenfolge existiert auch eine Inklusionskette der definierten Räumlichkeiten:
  $$
  \GConv \supset \KConv \supset \Lim \supset \PsTop \supset \PrTop \supset \TPrTop.
  $$
\end{bem}

\begin{proof}
  Jeder topologische Raum ist per definitionem ein prätopologischer Raum.

  Jeder prätopologische Raum ist ein pseutopologischer Raum: 
  Ist nämlich $(X,q) \in |\PrTop|$, so gilt $(\F,x) \in q$ genau dann, wenn $\F \supset U_q(x)$. Setzen wir nun voraus, dass $(\U,x) \in q$ für alle Ultrafilter $\U \supset \F$ gilt, so folgt aus
  $$
  \U_q(x) \subset \bigcap \{ \U \colon \U \in \FU(\F)\} = \F,
  $$
  wobei $\FU(\F)$ die Menge der Oberultrafiter von $\F$ bezeichne, die Behauptung durch Anwendung von C2).

  Jeder pseutotopologische Raum ist ein Limesraum: 
  Angenommen C4) sei nicht erfüllt für einen Limesraum $(X,q)$, so existieren Filter $\F,\G \in \FF(X)$ mit $(\F,x) \in q$ und $(\G,x) \in q$ aber $(\F \cap \G,x) \not\in q$.
  Folglich besitzt $(\F \cap \G, x)$ nach C5) ein Oberultrafilterfilter $(\U,x) \not\in q$.
  Insbesondere gilt nach C2) $\U \not\supset \F$ und $\U \not\supset \G$, es existieren also $F \in \F$ und $G \in \G$ mit $F,G \not \in \U$.
  Da $\U$ jedoch ein Oberfiter von $\F \cap \G$ ist, enthält er $F \cup G$ und aufgrund der Ultrafiltereigenschaft $F$ oder $G$ im Widerspruch zu $F, G \not\in \U$. 

  Jeder Limesraum ist ein Kent Konvergenzraum: Dies folgt sofort aus C1).

  Dass jeder Kent Konvergenzraum ein verallgemeinerter Konvergenzraum ist, ist wie bei allen anderen Konvergenzstrukturen Teil der Definition.
\end{proof}

\begin{prop}
  $\KConv$ ist bireflektives und bicoreflektives Unterkonstrukt von $\GConv$.
\end{prop}

\begin{proof}
  Sei $(X,q)$ ein verallgemeinerter Konverenzraum.
  Es lassen sich wie folgt zwei Kent Konvergenzstrukturen $q_r,q_c$ auf $X$ definieren:
  \begin{align*}
    &(\F,x) \in q_r \iff \exists (\G, x) \in q \text{, sodass } \G \cap \dot x \subset \F, \\
    &(\F,x) \in q_c \iff (\F, x) \in q \text{ und } (\F \cap \dot x) \in q.
  \end{align*}
  Es ist klar, dass es sich bei beiden Konvergenzstrukturen, um Kent Konvergenzstrukturen handelt.
  Aus der Konstruktion ergibt sich zudem $q \subset q_r$ und damit $1_X \colon (X,q) \to (X, q_r) \in \Mor_{\GConv}$ sowie $q_s \subset q$ und folglich $1_X \colon (X, q_s) \to (X,q) \in \Mor_{\GConv}$.
  Ist nun $(Y,u) \in |\KConv|$ und $f \colon (X, q) \to (Y,u)$ ein $\GConv$-Morphismus so ist durch $g \coloneqq 1_X^{-1} \circ f$ der gesuchte $\KConv$-Morphismus für die Faktorisierung gegeben.
  Damit ist $1_X$ die gesuchte Bireflektion.
  Analog beweist man den Fall für die Struktur $q_s$, für die $1_X$ zu einer Koreflektion wird.
\end{proof}

\begin{prop}[\cite{preuss}, 2.3.1.5]
  \label{prop:inklusion}
  Jedes der Konstrukte der Inklusionskette
  $$
  \KConv \supset \Lim \supset \PsTop \supset \PrTop \supset \TPrTop.
  $$
  ist ein bireflektives, volles und unter Isomorphie abgeschlossenes Unterkonstrukt der vorangehenden.
\end{prop}

Es mag die Frage aufgekommen sein, weshab wir die Objekte der Kategorie $\TPrTop$ als topologische Räume bezeichnet haben.
Zwischen besagter Kategorie $\TPrTop$ und $\Top$ besteht eine besondere Art der Isomorphie, die wie sollte es auch anders sein, sich einen Namen gemacht hat:
\begin{defn}
  \begin{enumerate}[a)]
    \item Seien $\A$ und $\B$ Kategorien.
      Dann nennen wir einen Funktor $\F \colon \A \to \B$ einen \emph{Isomorphismus}, falls ein Funktor $\G \colon \B \to \A$ existiert, sodass $\G \circ \F = \I_\A$ und $\F \circ \G = \I_\B$ gilt.
    \item Seien $\A$ und $\B$ Konstrukte.
      Sind $\Ha \colon \A \to \Set$ und $\K \colon \B \to \Set$ die Vergissfunktoren, so nennen wir einen Funktor $\A \to \B$ \emph{konkret}, falls $\K \circ \F = \Ha$ gilt.
    \item Einen konkreten Funktor $\F \colon \A \to \B$, welcher zudem ein Isomorphismus ist, bezeichnen wir als \emph{konkreten Isomorphismus}.
    \item Wir nennen zwei Kategorien (zwei Konstrukte) $\A$ und $\B$ \emph{isomorph} (\emph{konkret isomorph}), vorausgesetzt es existiert ein Isomorphismus (konkreter Isomorphismus) $\F \colon \A \to \B$.
  \end{enumerate}
\end{defn}

\begin{prop}[\cite{preuss}, 2.3.18]
  $\Top$ und $\TPrTop$ sind konkret Isomorph.
\end{prop}


\subsection{Uniforme Konvergenzstrukturen}

In diesem Abschnitt besprechen wir das zweite für uns interessante topologische Konstrukt mitsamt interessanter Unterkonstrukte.

\begin{defn}[SUConv und Nachfahren]
  Die Kategorie $\SUConv$ der verallgemeinerten Konvergenzräume mit gleichmäßig stetigen Abbildungen setzt sich wie folgt zusammen:
  \begin{enumerate}[a)]
    \item Ein \emph{semiuniformer Konvergenzraum} is ein Paar $(X, \J_X)$, wobei $X$ eine Menge und $\J_X \subset F(X \times X)$ die Menge der \emph{uniformen Filter} ist mit folgenden Eigenschaften:
      \begin{enumerate}[UC1)]
         \item $(\dot x \times \dot x) \in \J_X$ für alle $x \in X$.
         \item $\G \in \J_X$, falls $\F \in \J_X$ und $\F \subset \G$.
         \item Aus $\F \in \J_X$ folgt $\F^{-1} = \{F^{-1} \colon F \in \F\} \in \J_X$.
      \end{enumerate}
    \item Eine Abbildung $f \colon (X, \J_X) \to (Y, \J_Y)$ zwischen semiuniformen Konvergenzräumen heißt \emph{gleichmäßig stetig}, falls $(f \times f)(\J_X) \subset \J_Y$ gilt.
  \end{enumerate}
  Ein semiuniformer Konvergenzraum heißt 
\begin{enumerate}[a)]
  \item[c)] \emph{semiuniformer Limesraum}, falls die folgende Bedingung erfüllt ist:
    \begin{enumerate}[UC1)]
      \item[UC4)] $F \in \J_X$ und $\G \in \J_X$ implizieren $\F \cap \G \in \J_X$.
    \end{enumerate}
  \item[d)] \emph{uniformer Limesraum}, falls die folgende Bedingung erfüllt ist:
    \begin{enumerate}[UC1)]
      \item[UC5)] $\F \in \J_X$ und $\G \in \J_X$ implizieren $\F \circ \G \in \J_X$.
    \end{enumerate}
  \end{enumerate}
  Ein uniformer Limesraum $(X, \J_X)$ heißt
  \begin{enumerate}[a)]
    \item[e)] \emph{Haupt-uniformer Limesraum} falls eine nichtleere Teilmenge $\F$ von $\Pot(X \times X)$ existiert, welche die endliche Durchschnitsseigenschaft besitzt und gegenüber Obermengenbildung abgeschlossen ist und $[F] \coloneqq \{\G \in F(X\times X) \colon \G \supset \F \}= \J_X$ erfüllt.
  \end{enumerate}
\end{defn}

Die eben definierten Klassen definieren volle und unter Isomorphie abgeschlossene Unterkonstrukte von $\SUConv$, welche wir im Folgenden mit $\SULim$, $\ULim$ und $\PrULim$ bezeichnen werden.
Auch für diese Unterkonstrukte existiert eine Inklusionsbeziehung

\begin{prop}
  Jedes der Konstrukte der Inklusionskette
  $$
  \SUConv \supset \SULim \supset \ULim \supset \PrULim
  $$
  ist ein bireflektives, volles und unter Isomorphie abgeschlossenes Unterkonstrukt der vorangehenden.
\end{prop}


\subsection{Das Bindeglied zwischen Konvergenzstrukturen und uniformen Konvergenzstrukturen}

Haben wir in den beiden vorangehenden Sektionen Konvergenzstrukturen und uniforme Konvergenzstrukturen getrennt betrachtet, so kümmern wir uns nun darum die Verbindung zwischen beiden Strukturen herzustellen.
Zuvor will jedoch unser Vokabelheft gefüttert werden.

\begin{defn}
  Die Kategorie $\Fil$ der Filterräume und Cauchy stetigen Abbildungen setzt sich wie folgt zusammen:
  \begin{enumerate}[a)]
    \item Ein \emph{Filterraum} ist ein Paar $(X, \gamma)$, wobei $X$ eine Menge und $\gamma$ eine Menge von Filtern ist, sodass die Folgenden Bedingungen erfüllt sind.
      \begin{enumerate}[F1)]
        \item $\dot x \in \gamma$ für alle $x \in X$.
        \item $\G \in \gamma$, falls $\F \in \gamma$ und $\F \subset \G$.
      \end{enumerate}
      Ist $(X,\gamma)$ ein Filterraum, so wollen wir die Elemente von $\gamma$ \emph{Cauchy-Filter} nennen.
    \item Eine Abbildung $f \colon (X, \gamma) \to (X', \gamma')$ zwischen Filterräumen heißt \emph{Cauchy-stetig}, falls $f(\F) \in \gamma'$ gilt, für alle $\F \in \gamma$.
  \end{enumerate}
\end{defn}

\begin{defn}
  Sei $(X, \J_X)$ ein semiuniformer Konvergenzraum.
  Wir nennen einen Filter auf $X$ einen $\J_X$-\emph{Cauchy-Filter}, wenn $\F \times \F \in \J_X$ gilt.
\end{defn}

\begin{defn}
  Wir nennen einen semiuniformen Konvergenzraum $(X, \J_X)$ $\Fil$-bestimmt, falls $\J_X = \J_{\gamma_{J_X}}$ gilt, also $\J_X$ von allen $\J_X$-Cauchy-Filtern
\end{defn}

Wer bis hierhin gelesen hat, ahnt bereits, was kommen wird: Unser Bindeglied zwischen Konvergenzstrukturen und uniformen Konvergenzstrukturen werden die Filterräume sein. Das ist schon fast richtig. \emph{Fast} im Sinne von \emph{bis auf konkrete Isomorphie}.


\begin{prop}[\cite{preuss}, 2.3.3.5]
  Ist $\FilDSUConv$ das Konstrukt aller $\Fil$-bestimmten semiuniformen Konvergenzräume (mit gleichmäßig stetigen Abbildungen), dann ist $\Fil$ konkret isomorph zu $\FilDSUConv$.
\end{prop}

Nun geht es ans Eingemachte: Wir stellen die erste Verbindung zu uns bekannten Strukturen her.

\begin{prop}
  $\FilDSUConv$ ist ein bireflektives und bikoreflektives, volles und unter Isomorphie abgeschlossenes Unterkonstrukt von $\SUConv$.
\end{prop}

Nun steht nur noch offen eine Verbindung zu Konvergenzstrukturen herzustellen.
Dies funktioniert jedoch nicht unmittelbar. 
Um uns im bevorstehenden Terrain weiter bewegen zu können notieren wir in unserem Vokabelheft:

\begin{defn}
  Ein Filterraum $(X, \gamma)$ heißt \emph{vollständig}, falls für alle $\F \in \gamma$ ein $x \in X$ existiert mit $\F \cap \dot x \in \gamma$.
\end{defn}

Damit definieren wir uns nun ein neues Objekt, welches wir in Zusammenhang mit Konvergenzstrukturen bringen werden.

\begin{prop}
  Das Konstrukt $\CFil$ der vollständigen Filterräume (und Cauchy-stetigen Abbildungen) ist ein volles und unter Isomorphie abgeschlossenes, bicoreflektives Unterkonstrukt von $\Fil$.
\end{prop}

Die Konvergenzstrukturen, die sich mit $\CFil$ in Verbindung bringen lassen, besitzen eine Symmetrieeigenschaft, welche den Zoo der uns bekannten Konvergenzstrukturen mit weiteren bisher unerkannten Arten ausstattet.

\begin{defn}
  Ein verallgemeinerter Konvergenzraum $(X,q)$ heißt \emph{symmetrisch}, falls folgende Bedingung erfüllt ist:
  \begin{enumerate}[S1)]
    \item[S)] $(F,x) \in q$ und $y \in \bigcap_{F \in \F}$ implizieren $(F, y) \in q$.
  \end{enumerate}
\end{defn}

Nun zu dem Kandidaten, welcher Filterräume und Konvergenzräume verbindet.

\begin{prop}
  \begin{enumerate}[1)]
    \item \begin{enumerate}[a)]  
        \item Sei $(X, \gamma)$ ein Filterraum. Dann wird eine symmetrische Kent Konvergenzstruktur $q_\gamma$ auf $X$ definiert durch
          $$
          (\F, x) \in q_\gamma \iff \F \cap \dot x \in \gamma
          $$
        \item Ist $f \colon (X, \gamma) \to (X', \gamma')$ eine Cauchy-stetige Abbildung zwischen Filterräumen, so ist $f \colon (X, q_\gamma)$ stetig.
      \end{enumerate}
    \item \begin{enumerate}[a)]
        \item Sei $(X, q)$ ein Kent Konvergenzraum. Dann wird eine vollständige $\Fil$-Struktur $\gamma_q$ auf $X$ definiert durch
          $$
          \gamma_q = \{ \F \in \FF(X) \colon \exists x \in X \colon (\F, x ) \in q \}.
          $$
        \item Ist $f \colon (X, q) \to (X', q')$ eine stetige Abbildung zwischen Kent Konvergenzräumen, dann ist $f \colon (X, \gamma_q) \to (X', \gamma_{g'})$ Cauchy-stetig.
      \end{enumerate}
    \item Das Konstrukt $\CFil$ ist konkret isomorph zum Konstrukt $\KConv_\Sy$ der symmetrischen Kent Konvergenzräume (und stetigen Abbildungen).
  \end{enumerate}
\end{prop}

Was hate es nun mit symmetrischen Konvergenzräumen auf sich und wie übertragen sich die Eigenschaften aus Proposition \ref{prop:inklusion} auf ihre symmetrischen Verwandten?

\begin{prop}
  \begin{enumerate}[1)]
    \item Sei $\A$ ein topologisches Konstrukt der Klasse Limesraum, pseudotopologischer Raum, prätopologischer Raum oder topologischer Raum.
      Dann ist das volle und unter Isomoprhie abgeschlossene Unterkonstrukt $\A_\S$ bestehend aus symmetrischen Objekten bireflektiv in $\A$.
    \item Jedes der Konstrukte der Inklusionskette
      $$
      \KConv_\Sy \supset \Lim_\Sy \supset \PsTop_\Sy \supset \PrTop_\Sy \supset \Top_\Sy
      $$
      ist ein bireflektives Unterkonstrukt der Vorgänger.
  \end{enumerate}
\end{prop}


Folgendes Diagramm liefert eine Zusammenfassung der besprochenen Sachverhalte:
$$
\begin{tikzcd}[column sep=5.5em,row sep=2.5em]
  \SUConv              &        & \GConv \\
  \SULim \arrow[u,"r"] &  \Fil\arrow[lu,"rc"]  & \KConv   \arrow[u,"rc"] \\
  \ULim  \arrow[u,"r"] &        & \KConv_\Sy \arrow[lu,"c"]    \arrow[u,"r"]    \\
  \Unif  \arrow[u,"r"] &        & \Lim_\Sy   \arrow[u,"r"]    \\
                       &        & \PsTop_\Sy \arrow[u,"r"]    \\
                       &        & \PrTop_\Sy \arrow[u,"r"]    \\
                       &        & \Top_\Sy   \arrow[u,"r"]    
\end{tikzcd}
$$


\bibliographystyle{alpha}
\bibliography{seminar17_fabian_gabel}
      
\end{document}
